\section{An absolutely convergent formula for the normalization}
\label{sec:normalization}

The spectral constant can be specified entirely by finite counting problems.
This section fixes every ingredient of the formula and proves convergence of
its indicated terms. The construction is independent of a domain and its
marked boundary points.

Choose the even smooth cutoff
\begin{equation}
 \label{eq:explicit-cutoff}
 q(s)=
 \begin{cases}
  1,& |s|\le1/2,\\[2pt]
  \displaystyle
  \frac{e^{-1/(1-|s|)}}
       {e^{-1/(1-|s|)}+e^{-1/(|s|-1/2)}},&1/2<|s|<1,\\[8pt]
  0,&|s|\ge1.
 \end{cases}
\end{equation}
For \(j=0,1,2,\ldots\), define a smooth function on the reciprocal torus by
\begin{equation}
 \label{eq:normalization-multiplier}
 F_j(p)=q(4|p|)\cos(4\arg p)
                    \bigl[q(2^jk)-q(2^{j+1}k)\bigr],\qquad p=(k,l),
\end{equation}
using its expression in the disk \(|p|<1/4\), and zero outside that disk.
Set \(F_j(0)=0\). Smoothness at the apparent angular singularity follows
because the bracket vanishes whenever \(|k|\le2^{-j-2}\).
Let \(d_{ju}\), \(u\in\Tri\), be its Fourier coefficients, with the convention
\begin{equation}
 \label{eq:normalization-fourier}
 F_j(p)=\sum_{u\in\Tri}d_{ju}e^{ip\cdot u},\qquad
 d_{ju}=\frac{A_{\Tri}}{(2\pi)^2}
        \int_{\widehat\Tri}F_j(p)e^{-ip\cdot u}\,dp,
 \qquad A_{\Tri}=\frac{\sqrt3}{2}.
\end{equation}
Here the integral uses Euclidean measure on any fundamental representative
of the torus, whose area is \((2\pi)^2/A_{\Tri}\). The coefficients are real
and satisfy \(d_{j,-u}=d_{ju}\), because \(F_j\) is real and even.

For \(j,s\ge0\), put
\begin{equation}
 \label{eq:normalization-radii}
 R_{js}=2^{(j+s+4)^2},\qquad
 \mathcal H_R=\{u\in\Tri:\dist_{\Tri}(0,u)\le R\}.
\end{equation}
Let \(\langle\cdot\rangle_R\) be uniform averaging over all pairs
\(B,W\in\{-1,+1\}^{\mathcal H_R}\) for which both colors do not change on the
same nearest-neighbor edge. There are no boundary pins. For
\(|u|\le\sqrt R\), choose any lattice path \(\gamma_{0u}\) in
\(\mathcal H_R\), oriented from \(0\) to \(u\), and set
\begin{equation}
 \label{eq:finite-band-coefficient}
 b_{js}=\sum_{\substack{u\in\Tri\\ |u|\le\sqrt{R_{js}}}}d_{ju}
  \left\langle
   \left(\sum_{xy\in\gamma_{0u}}
           \frac{W_xB_y-B_xW_y}{2}\right)^2
  \right\rangle_{R_{js}}.
\end{equation}
The curl identity for the spin encoding makes this independent of the paths.
The expression is a finite sum of prescribed Fourier coefficients and ratios
of finite integer counts.

\begin{proposition}[Finite-volume normalization series]
\label{prop:normalization-series}
The series of the indicated terms is absolutely convergent:
\begin{equation}
 \label{eq:normalization-absolute}
 \sum_{j\ge0}|b_{j0}|+
       \sum_{j,s\ge0}|b_{j,s+1}-b_{js}|<\infty.
\end{equation}
Its value satisfies
\begin{equation}
 \label{eq:finite-normalization}
 c_*=\frac1\pi\sum_{j\ge0}
       \left[b_{j0}+\sum_{s\ge0}(b_{j,s+1}-b_{js})\right]>0,
 \qquad
 \sigma=4\pi\sqrt{c_*}.
\end{equation}
Thus the displayed formula specifies a strictly positive universal
normalization using only \cref{eq:explicit-cutoff,eq:normalization-multiplier,eq:normalization-fourier,eq:normalization-radii,eq:finite-band-coefficient}.
\end{proposition}

\begin{proof}
Write
\[
 V(u)=\EE\bigl(H(u)-H(0)\bigr)^2,\qquad
 b_j=\sum_{u\in\Tri}d_{ju}V(u).
\]
The height rule gives \(V(u)\le K|u|^2\) for \(u\ne0\). For each fixed \(j\),
Fourier smoothness makes the series defining \(b_j\) absolutely convergent.
By \cref{eq:increment-spectral},
\[
 V(u)=2\int(1-\cos(p\cdot u))\,C(dp).
\]
Interchanging the sum and integral is legitimate even though \(C\) has
infinite mass near zero: the absolute integral is bounded by
\(\sum_u|d_{ju}|V(u)<\infty\). Since
\(\sum_u d_{ju}=F_j(0)=0\) and
\(\sum_u d_{ju}\cos(p\cdot u)=F_j(p)\), we obtain
\begin{equation}
 \label{eq:band-inversion}
 b_j=-2\int F_j(p)\,C(dp)=-2I_j.
\end{equation}
It follows from \cref{prop:angular-sum} that
\begin{equation}
 \label{eq:infinite-band-sum}
 \sum_{j\ge0}|b_j|<\infty,
 \qquad \sum_{j\ge0}b_j=\pi c_*.
\end{equation}

We next bound the finite-volume error uniformly in both indices. Differentiating
\cref{eq:normalization-multiplier} a fixed number of times gives
\begin{equation}
 \label{eq:band-fourier-decay}
 |d_{ju}|\le K\,2^{Aj}(1+|u|)^{-8}
\end{equation}
for fixed finite constants \(K,A\). Indeed, on the support of the horizontal
bracket we have \(|p|\ge|k|\ge2^{-j-2}\); derivatives of the angular factor
and of the rescaled cutoff therefore cost only a fixed power of \(2^j\).
Integration by parts in a basis of torus coordinates gives the displayed
decay.

Let \(V_R(u)\) denote the finite-hexagon expectation in
\cref{eq:finite-band-coefficient}. Reduce the exponent in
\cref{prop:plane-mixing}, if necessary, to \(0<\alpha\le1\). For
\(0<|u|\le\sqrt R\), evaluate the increment along a shortest path contained
in a ball of radius \(K|u|\). This is permitted by path independence. Both
its finite and infinite versions have square at most \(K|u|^2\). Coupling
inside this ball gives
\begin{equation}
 \label{eq:finite-variance-error}
 |V_R(u)-V(u)|\le K|u|^2\left(\frac{|u|}{R}\right)^\alpha.
\end{equation}
The distance to the free outer boundary is comparable to \(R\), uniformly in
the stated range. Combining
\cref{eq:band-fourier-decay,eq:finite-variance-error}, the error in the retained
part of the coefficient is at most
\[
 K2^{Aj}R^{-\alpha}
       \sum_{u\in\Tri}(1+|u|)^{-8}|u|^{2+\alpha}
   \le K'2^{Aj}R^{-\alpha}.
\]
The omitted tail is at most
\[
 K2^{Aj}\sum_{|u|>\sqrt R}(1+|u|)^{-8}|u|^2
       \le K'2^{Aj}R^{-2}.
\]
Consequently, for some fixed \(\beta>0\),
\begin{equation}
 \label{eq:finite-band-error}
 |b_{js}-b_j|\le K2^{Aj}R_{js}^{-\beta}.
\end{equation}
All constants here are independent of \(j,s\).

Put \(e_{js}=b_{js}-b_j\). The particular radii in
\cref{eq:normalization-radii} make
\(\sum_{j,s}|e_{js}|<\infty\), since a negative quadratic exponent dominates
\(Aj\). Moreover,
\[
 |b_{j,s+1}-b_{js}|\le|e_{j,s+1}|+|e_{js}|,
 \qquad |b_{j0}|\le|b_j|+|e_{j0}|.
\]
Together with \cref{eq:infinite-band-sum}, this proves
\cref{eq:normalization-absolute}. For every \(j\), telescoping gives
\(b_{j0}+\sum_s(b_{j,s+1}-b_{js})=b_j\), and then
\cref{eq:infinite-band-sum} gives the first identity in
\cref{eq:finite-normalization}. Its positivity was proved in
\cref{thm:plane-covariance}. Finally, the covariance coefficient for \(H\) is
\(v=(2\pi)^2c_*\). The centered odd height is twice the centered integer
height, so its covariance coefficient is \(4v\), whose positive square root
is \(4\pi\sqrt{c_*}\).
\end{proof}

\begin{remark}[The fundamental-cell convention]
\label{rem:cell-area}
There is no further triangular-cell factor in
\cref{eq:finite-normalization}. If \(\lambda_u\) is the unit-mesh nodal hat,
then \(\int\lambda_u=A_{\Tri}\). Thus the physical interpolation weight is
\[
 a_{\delta,u}(f)
   =\delta^2\int_{\R^2}\lambda_0(y)f(\delta u+\delta y)\,dy.
\]
The hat mass \(A_{\Tri}\delta^2\) cancels the site density
\((A_{\Tri}\delta^2)^{-1}\) in the limit of the transform at
\(p=\delta\zeta\). Its limit is \(\widehat f(-\zeta)\), with no prefactor.
The reciprocal-cell factor in \cref{eq:normalization-fourier} is already part
of the Fourier coefficients used in \cref{eq:band-inversion}. The only
remaining comparison is between
\(c_*\int|\widehat f(\zeta)|^2|\zeta|^{-2}\,d\zeta\) and the conventional
inverse-Laplacian covariance
\((2\pi)^{-2}\int|\widehat f(\zeta)|^2|\zeta|^{-2}\,d\zeta\).
This accounts for \(v=(2\pi)^2c_*\).
\end{remark}

With the normalization fixed, we return to the reflection-null identity of
\cref{cor:reflection-null}. In \cref{sec:pins} we transfer that identity to
laws conditioned on the boundary pins.
